\begin{lemma}[Homogeneous triples]
For every infinite $X\subseteq\mathbb N$ and every two-coloring of the
three-element subsets of $X$, there are an infinite $Y\subseteq X$ and a
color $b$ such that every three-element subset of $Y$ has color $b$.
\end{lemma}
\begin{proof}
Let $F=\mathsf{CardinalityFront}(3,X)$, as defined in
\noderef{CardinalityFront}.  By
\noderef{CardinalityFrontIsFront}, $F$ is a front on $X$.  Define

\[
S=\{s\in F\mid c(s)=1\}.
\]

We first check the hypothesis needed for Nash--Williams.  If $s\in S$,
then the set-builder definition of $S$ gives $s\in F$ directly.  Hence
$S\subseteq F$.  Applying \noderef{NashWilliams} to $F$, $X$, and $S$
therefore gives a set $F'$ and an infinite set $Y$ such that

\[
\mathsf{Front}(F',Y),\qquad F'\subseteq F,
\qquad F'\subseteq S\ \text{or}\ F'\cap S=\emptyset.
\tag{1}
\]

We next show that $F'$ is not the trivial front.  The infinite-base and
density clauses of \noderef{Front} give $Y$ infinite and, after applying
density to the infinite subset $Y$ of this front's base parameter, a set
$t\in F'$ with $\mathsf{ProperPrefixSet}(t,Y)$.  Since $F'\subseteq F$,
the definition of \noderef{CardinalityFront} gives $t\subseteq X$ and
$|t|=3$.  In particular $t\ne\emptyset$.  Thus $F'\ne\{\emptyset\}$:
if equality held, the element $t\in F'$ just obtained would be empty,
contradicting $|t|=3$.

The base disjunction in \noderef{Front} now forces

\[
\mathsf{FrontBase}(F')=Y.
\tag{2}
\]

Consequently $Y\subseteq X$.  Indeed, if $y\in Y$, equation (2) and
the definition of \noderef{FrontBase} give an $s\in F'$ with $y\in s$.
The inclusion $F'\subseteq F$, together with
\noderef{CardinalityFront}, says that every element of $s$ belongs to
$X$, so $y\in X$.

We claim that every three-element subset of $Y$ belongs to $F'$.  Let
$s\subseteq Y$ have three elements.  Write its elements in increasing
order as

\[
a<b<c,\qquad s=\{a,b,c\}.
\]

The tail

\[
Y_{>c}=\{y\in Y\mid c<y\}
\]

is infinite.  For if it were finite, then
$Y\cap\{0,1,\ldots,c\}$ would also be finite and

\[
Y=(Y\cap\{0,1,\ldots,c\})\cup Y_{>c}
\]

would be finite, contrary to the infinitude of $Y$.  Put

\[
Z=\{a,b,c\}\cup Y_{>c}.
\]

Then $Z\subseteq Y$ is infinite, and its only elements at most $c$ are
$a,b,c$.  Apply the density clause of \noderef{Front} to $Z$.  We get
$t\in F'$ and \mathsf{ProperPrefixSet}(t,Z)$.  Unpacking
\noderef{ProperPrefixSet}, choose $n\in Z$ such that

\[
\forall k,\qquad k\in t\quad\Longleftrightarrow\quad
k\in Z\ \text{and}\ k<n.
\tag{3}
\]

Again $t\in F'\subseteq F$, so \noderef{CardinalityFront} gives
$|t|=3$.  Since $n\in Z$, there are four possibilities.  If $n=a$,
the right side of (3) is empty; if $n=b$, it is $\{a\}$; and if
$n=c$, it is $\{a,b\}$.  Each of these contradicts $|t|=3$.  Thus
$c<n$.  In this case (3) shows that $a,b,c\in t$, so
$\{a,b,c\}\subseteq t$.  If some $y\in Y_{>c}$ also satisfied $y<n$,
then (3) would put the four distinct elements $a,b,c,y$ in $t$,
contradicting $|t|=3$.  Hence no element of $Y_{>c}$ is below $n$.
Equation (3) then says precisely that $t=\{a,b,c\}=s$.  Therefore
$s\in F'$, proving the claim.

Finally consider the two alternatives in (1).  If $F'\subseteq S$,
then every three-element $s\subseteq Y$ lies in $S$, so the definition
of $S$ gives $c(s)=1$; choose $b=1$.  If
$F'\cap S=\emptyset$, then every such $s$ lies in $F'$ but not in $S$.
Since $s\in F$ and $S=\{u\in F\mid c(u)=1\}$, this means
$c(s)\ne1$.  A Boolean value is either $0$ or $1$, so $c(s)=0$; choose
$b=0$.  In both cases $Y\subseteq X$, $Y$ is infinite, and every
three-element subset of $Y$ has color $b$, as required.
\end{proof}
